% Fragmento público generado desde propietarios íntegros.
% Residencia: 52.1.
% Fuente: ../incoming/radion_monografia_20260827/manuscrito/sections/01_resumen_y_guia.tex:1-128
\paragraph{Mathematical Synthesis of the Object and Its Causal Chain}

This monograph reconstructs the radian from the generative chain
\[
\APP\longrightarrow\TRIT\longrightarrow\TPK
\longrightarrow X_{\mathrm{enr}}
\longrightarrow\text{discrete structure of the continuum}
\longrightarrow\MD,
\]
without taking as input either the metrological value of \(180/\pi\) or the small angular residue to be explained. The result is an enriched unit, the \emph{radion}, defined as the oriented unit that traverses one radian of phase and sweeps exactly an action \(\hret\) in the positive ellipse associated with the same state.

The governing closure is
\BigRadionEquation
with
\[
S_E=2\pih\left(\frac5{36}+\frac{25}{9}\alphah\right),
\qquad
\Delta_4=
\frac{\pih-e_{\HMT}\log\pih}{270}
\left(1+\frac{\pih}{729}\right),
\]
\[
\rho_{\mathrm{tw}}=\frac{2\cdot90}{729}=\frac{20}{81},
\qquad
\epsone=\frac{20}{81}\Delta_4-(S_E-1).
\]
APP division proves \(1<S_E<2\), so that the quotient is exactly
one and the remainder is \(D_E=S_E-1\). The equality
\[
1=S_E-\frac{20}{81}\Delta_4+\epsone
\]
is then an internal identity. The angular chart produces
\[
\epsR^\circ
=5.13830167585752820716851646226459474899085744\ldots
\times10^{-8}\,{}^\circ.
\]

The text also demonstrates that the action ellipse
\[
Q=a_S\cos\theta,\qquad P=b_S\sin\theta,\qquad
a_Sb_S=2\hret
\]
satisfies
\[
\operatorname{Area}(0,\theta)=\hret\theta.
\]
Therefore, a radion sweeps \(\hret\) and the turn \(2\pi\) encloses
\(h_{\mathrm{ret}}=2\pi\hret\). The same ellipse has a finite symplectic
boundary and an interior of infinite area when equipped with the
Hilbert--Beltrami--Klein metric. This is the exact form of the physical thesis guiding
the work: the visible surface closes the action; the interior depth
preserves a development that cannot be exhausted by the planar projection.

The exposition maintains typed four realizations of the same state:

\begin{enumerate}[label=\textbf{\arabic*.}]
\item the circle \(S^1\), which publishes the phase;
\item the positive ellipse, which publishes the action and the anisotropy;
\item the interior of Klein, that miof the depth by means of reason doble;
\item the nonadic helix, which preserves return, memory, and contraction.
\end{enumerate}

They are not four superposed geometries without a criterion. They are four readers with distinct domains, codomains, and conserved data, applied to a single enriched state. The sector \(90/120\), the dodecaphase, the radion carry, the Lambert tail, the Barbero kernel, and the global towers are likewise presented as related but noninterchangeable objects.

\begin{thesisbox}
The radian is the oriented unit that sweeps \(\hret\) in the action ellipse.
The circle publishes its phase, Klein measures its depth, the helix preserves its
memory, and the coinductive carry exactly splices its direct and
torsional sections.
\end{thesisbox}

\paragraph{Expository order and reading dependencies}

The order is not ornamental. Parts I and II produce the object and its
closure from APP--TRIT--TPK. Only thereafter are realizations introduced in
the languages of Hilbert, LC, Maxwell, Minkowski, Hodge, Lorentz,
modularity, and orbits. Each chapter distinguishes three levels:

\begin{description}[style=nextline,leftmargin=4.2cm,labelwidth=3.9cm]
\item[Internal result.] Equality or construction proved in the
HMT domain with its declared primitives.
\item[Consolidated formalization.] A new map articulating already
existing results without using the target as a selector.
\item[Physical interface.] Subsequent dimensional translation that recognizes a
conventional structure without making it the origin of the selection.
\end{description}

The \emph{status and falsifier} boxes indicate which concrete remark would invalidate each strong promotion. This discipline does not weaken the thesis; it makes the thesis auditable. The purpose of the extended narrative is to ensure that no word such as \emph{circle}, \emph{ellipse}, \emph{torsion}, \emph{charge}, or \emph{memory} conceals a change of type.

\paragraph{Pivotal results and proof scope}

\begin{enumerate}[label=\textbf{C\arabic*.}]
\item Closure of the radion independent of its final value through two publications of the same state and
APP subtraction with carry.
\item Derivation of the coefficient \(20/81\) as two sheets of the active sector
\(90\), normalized by the fiber \(3^6=729\).
\item Typed equalizer between the seam \(\Re s=1/2\), the hundred-mark APP
face, and the dodecaphasic phase \(\theta_2=50^\circ\).
\item Radian area theorem: \(\operatorname{Area}(\theta)=\hret\theta\).
\item Exact focal study of the ellipse and of its Euclidean,
hyperbolic, symplectic, and modular invariants.
\item Boundary--interior theorem: finite total action \(h\) and Klein interior
of infinite hyperbolic area.
\item Exact realization through minimal covariance and an LC oscillator, followed by
the electromagnetic publication \(90/120\).
\item Atlas of non-identifications for all objects \(90/120\), including
the separation among \(\epsR\), the spectral tail, and the Barbero kernel \(12:-1\).
\item Structural characterization of \(\pi\) in the radion chart, distinct
from the primary coinductive generator.
\item Ontological synthesis: dimension as the minimal rank of memories that a
phase projection can no longer preserve.
\end{enumerate}
% Fuente: ../incoming/radion_monografia_20260827/manuscrito/sections/01b_mapa_resultados.tex:1-133
\paragraph{Correspondence between results, causal types, and causal residences}

The following correspondence orders the governing assertions and their proofs,
while keeping visible the hierarchy that prevents
physical richness from erasing the mathematical closure supporting it.

% HMT_RESULT_BEGIN:RDN-01-GENEALOGIA:theorem
\begin{exactbox}[title={APP--TRIT--TPK genealogy of the radion}]
The radion arises from the effective composition
\(\APP\to\TRIT\to\TPK\to X_{\mathrm{enr}}\to\MD\); it is not obtained
by retroactively redefining the conventional radian.
\end{exactbox}
% HMT_RESULT_END:RDN-01-GENEALOGIA:theorem

% HMT_RESULT_BEGIN:RDN-02-CIERRE:theorem
\begin{exactbox}[title={Exact closure and prospective independence}]
APP carry internally produces
\[
1=S_E-\frac{20}{81}\Delta_4+\varepsilon_R^{(1)},
\qquad
\varepsilon_R^{(1)}=\frac{20}{81}\Delta_4-(S_E-1),
\]
and the angular chart transports this same identity to
\(180^\circ/\pi_{\!\HMT}\) without introducing the target residue.
\end{exactbox}
% HMT_RESULT_END:RDN-02-CIERRE:theorem

% HMT_RESULT_BEGIN:RDN-03-BISAGRA-50:theorem
\begin{exactbox}[title={Critical coordinate of \(50^\circ\)}]
The critical coordinate \(\Re s=1/2\), published on the APP face of one hundred
marks and in the dodecaphase wheel, selects the phase
\(\theta_2=50^\circ\). It is an equality of typed readers of the same state.
\end{exactbox}
% HMT_RESULT_END:RDN-03-BISAGRA-50:theorem

% HMT_RESULT_BEGIN:RDN-04-COEFICIENTE-2081:theorem
\begin{exactbox}[title={Nonadic Bidirectional Density \(20/81\)}]
The torsional coefficient is not fitted on the lattice \(1/81\): it is constructed as
\(\rho_{\mathrm{tw}}=2N_6/3^6=2\cdot90/729=20/81\), preserving the sheet,
orientation, and normalization of the nonadic fibre.
\end{exactbox}
% HMT_RESULT_END:RDN-04-COEFICIENTE-2081:theorem

% HMT_RESULT_BEGIN:RDN-05-AREA-ACCION:theorem
\begin{exactbox}[title={Areal law of the radion}]
For the positive ellipse
\(Q=a_S\cos\theta\), \(P=b_S\sin\theta\),
\(a_Sb_S=2\hbar_{\mathrm{ret}}\), the oriented area swept is
\(\mathcal A(\theta)=\hbar_{\mathrm{ret}}\theta\). One radion sweeps exactly
\(\hbar_{\mathrm{ret}}\), and the revolution \(2\pi\) encloses
\(h_{\mathrm{ret}}=2\pi\hbar_{\mathrm{ret}}\).
\end{exactbox}
% HMT_RESULT_END:RDN-05-AREA-ACCION:theorem

% HMT_RESULT_BEGIN:RDN-06-INVARIANTES-FOCALES:theorem
\begin{exactbox}[title={Focal invariants of the positive operator}]
The published pair \(A\pm C^\ast\) fixes the semiaxes, eccentricity,
focal separation, and symplectic compression parameter. In particular,
\((d_1/d_0)^2=C^\ast/2\), a dimensionless identity separating shape from scale.
\end{exactbox}
% HMT_RESULT_END:RDN-06-INVARIANTES-FOCALES:theorem

% HMT_RESULT_BEGIN:RDN-07-BORDE-INTERIOR:theorem
\begin{exactbox}[title={Symplectic finiteness of the boundary and divergence of the interior hyperbolic area}]
The symplectic boundary of the ellipse encloses the finite action
\(h_{\mathrm{ret}}\), whereas the Hilbert--Beltrami--Klein area of the interior
diverges upon reaching the boundary. Visible phase and depth are not the same
measure.
\end{exactbox}
% HMT_RESULT_END:RDN-07-BORDE-INTERIOR:theorem

% HMT_RESULT_BEGIN:RDN-08-HELICE-MEMORIA:theorem
\begin{exactbox}[title={Nonadic Helical Lift with Phase Return and Memory Advance}]
The nonadic monodromy returns the phase without resetting the state. The helix
lifts the circle while preserving carry and memory; the closure \(2\pi/4\pi\)
distinguishes projective return, reversal of orientation, and complete return.
\end{exactbox}
% HMT_RESULT_END:RDN-08-HELICE-MEMORIA:theorem

% HMT_RESULT_BEGIN:RDN-09-OSCILADOR-EM:development
\begin{narrativebox}[title={Oscillatory and constitutive electromagnetic realization}]
The same anisotropy is realized through a minimal covariance and an LC oscillator:
one coordinate stores the electric--elastic aspect and its conjugate the
magnetic--rotational aspect. Frequency and impedance remain separately typed.
\end{narrativebox}
% HMT_RESULT_END:RDN-09-OSCILADOR-EM:development

% HMT_RESULT_BEGIN:RDN-10-TIPOS-90120:theorem
\begin{exactbox}[title={Typed separation of the operators \(90/120\)}]
Dodecaphasic extractor, radion carry, Lambert tail, hexad--octad
functional, Barbero kernel, and global tower semigroup are distinct
objects. The weight \(12:-1\) comes from the twelve sectors of the
dodecaphasic fundamental chain and its unique oriented seam. The coefficient of
\((1-q)^{-1}\), equal to \(1/24\), and the minimal integer counterterm subsequently verify
the pair produced.
\end{exactbox}
% HMT_RESULT_END:RDN-10-TIPOS-90120:theorem

% HMT_RESULT_BEGIN:RDN-11-GEOMETRIAS-CAMPO:development
\begin{narrativebox}[title={Complex, Lorentzian, and Hodge-dual realizations}]
The elliptic, parabolic, and hyperbolic regimes of TRIT are published as
rotation, threshold, and hyperbolic transformation. The Hodge star, the constitutive relations, and the Lorentz
force act subsequently on the already oriented state; they do not select its
constants.
\end{narrativebox}
% HMT_RESULT_END:RDN-11-GEOMETRIAS-CAMPO:development

% HMT_RESULT_BEGIN:RDN-12-MODULAR-ORBITAL:development
\begin{narrativebox}[title={Modular, orbital, and arithmetic boundary characters}]
The modular involution, elliptic orbits, advance--delay, and
3--6--9 arithmetic recognize distinct memories of the same state. The modular function and
vacancies are used as posterior readers, never as substitutes for the
APP--TRIT--TPK generator.
\end{narrativebox}
% HMT_RESULT_END:RDN-12-MODULAR-ORBITAL:development

% HMT_RESULT_BEGIN:RDN-13-SINTESIS-ONTOLOGICA:conclusion
\begin{thesisbox}[title={Operatorial synthesis of phase, action, depth, and memory}]
The circle is the phase projection; the ellipse is the action section; Klein
measures depth; the helix preserves history. The radion is the oriented unit
that keeps these publications coupled without erasing the residue that
distinguishes them.
\end{thesisbox}
% HMT_RESULT_END:RDN-13-SINTESIS-ONTOLOGICA:conclusion

% HMT_RESULT_BEGIN:RDN-14-CONTROL-ALCANCE:control
\begin{statusbox}[title={Delimitation of domains and causal continuity}] The complete mass laws, electron refinement, CKM/CP, and the exceptional corridor retain their own governing proofs. This monograph cites them only where they share an antecedent or chart; it neither transfers their specific selectors to the radion nor confuses their probative residences.
\end{statusbox}
% HMT_RESULT_END:RDN-14-CONTROL-ALCANCE:control
